\newpage
% =================================================
\section{Implementações em Estudo para Versões Futuras}
% =================================================

\begin{explicacao}[Escopo desta seção]
Nada aqui faz parte do modelo 3FN da Seção 4 ou do mapeamento físico Firestore da Seção 5, nem deve ser implementado nesta primeira versão. É o registro de uma proposta para quando reagentes gasosos entrarem no almoxarifado do LCQUI -- hoje nenhum reagente gasoso está cadastrado, por isso \texttt{GASOSO} foi removido de \textit{Resumo\_Reagente.estado\_fisico} (seção 4.17) nesta versão, em vez de deixado junto de campos pensados só para sólido/líquido.
\end{explicacao}

\subsection{Evoluções posteriores da autorização}
M9 fixa a V1 em RBAC com escopo/vínculo persistido e ownership, sem introduzir
ACL, ABAC ou capabilities. Versões futuras só poderão propor novo paradigma
depois de demonstrar requisito de domínio que a matriz M9 não expressa, com
migração, regras, backend, revogação, auditoria, idempotência e casos de
segurança próprios. Não se amplia privilégio por conveniência de interface ou
por Custom Claim. A implementação do contrato M9 nas Rules, Storage, backend,
App Check, IAM e testes é trabalho posterior à formalização executável M9.

\subsection{Por que \texttt{GASOSO} sozinho não basta}

Simplesmente devolver \texttt{GASOSO} a \textit{estado\_fisico} resolveria a classificação básica, mas um cilindro de gás carrega informação que \textbf{não} é estado físico e que não deve virar uma explosão de valores dentro do mesmo enum (\texttt{GASOSO\_COMPRIMIDO}, \texttt{GASOSO\_LIQUEFEITO}, \texttt{MISTURA\_GASOSA\_CALIBRACAO}, \ldots). São quatro dimensões independentes, que devem continuar independentes no banco:

\begin{center}
\begin{tabular}{p{0.28\textwidth}p{0.62\textwidth}}
\toprule
\textbf{Dimensão} & \textbf{Exemplos de valor} \\
\midrule
Estado físico & sólido, líquido, gasoso \\
Forma de armazenamento & gás comprimido, gás liquefeito, gás criogênico \\
Tipo de substância (já existe) & pura, mistura \\
Finalidade/aplicação & reagente comum, gás de calibração, gás de processo \\
\bottomrule
\end{tabular}
\end{center}

Um cilindro de CO\textsubscript{2} puro, por exemplo, é \textit{estado\_fisico} = gasoso, \textit{forma\_armazenamento} = liquefeito (o CO\textsubscript{2} é mantido líquido sob pressão dentro do cilindro, embora seja gasoso em condição ambiente), \textit{tipo\_substancia} = pura. Já uma mistura de calibração (ex.: 500\,ppm de CO\textsubscript{2} balanceado em N\textsubscript{2}) é \textit{estado\_fisico} = gasoso, \textit{forma\_armazenamento} = comprimido, \textit{tipo\_substancia} = mistura, \textit{finalidade} = calibração. Tratar ``mistura gasosa de calibração'' como uma categoria física própria misturaria uma propriedade química (o que a substância é) com uma propriedade comercial/de uso (para que ela serve) -- o mesmo tipo de acoplamento indevido que a decomposição do \textit{status} de \textit{Frasco\_Reagente} (seção 4.19) já evitou para outro campo.

\subsection{Proposta de campos (apenas ilustrativa)}

\begin{entidade}{Especificacao\_Reagente \textit{(campos adicionais propostos)}}
\campo{forma\_armazenamento}{ENUM}{NULL}
\subcampo{CONVENCIONAL, GAS\_COMPRIMIDO, GAS\_LIQUEFEITO, GAS\_CRIOGENICO}
\campo{finalidade}{ENUM}{NULL}
\subcampo{REAGENTE, CALIBRACAO, PROCESSO, OUTRO}
\end{entidade}

Ambos ficam em \textit{Especificacao\_Reagente}: são propriedades do produto comercial específico, não da substância química abstrata nem do item de catálogo. \texttt{NULL} para itens sólidos/líquidos, onde a dimensão não se aplica.

\subsection{Além do enum: medição de gases não é pesagem}

A Regra de Cálculo de Volume e Peso (seção 6.2.13) assume que o volume é sempre derivado de uma leitura de balança convertida por densidade. Para cilindros de gás, a prática de laboratório mais comum é ler a \textbf{pressão} no manômetro do próprio cilindro (ainda que pesar o cilindro também seja uma alternativa viável em alguns casos). Isso implica que \textit{Frasco\_Reagente} precisaria, no mínimo, de:

\begin{itemize}
    \item \textit{pressao\_nominal\_bar} -- pressão de enchimento de fábrica.
    \item \textit{pressao\_atual\_bar} -- última leitura do manômetro (substituindo/complementando \textit{peso\_no\_cadastrado}/\textit{medida\_usada} para este tipo de frasco).
    \item \textit{capacidade\_litros\_cntp} -- volume de gás equivalente em CNTP (Condições Normais de Temperatura e Pressão) quando o cilindro está cheio.
    \item \textit{tipo\_valvula} -- necessário para compatibilidade de equipamento na retirada.
\end{itemize}

O cálculo de consumo também muda: em vez de $V = \Delta m / \rho$, a estimativa correta usa a proporcionalidade entre pressão e quantidade de gás restante (lei dos gases ideais, $PV = nRT$, ou uma correlação real para gases não-ideais em alta pressão), não uma conversão por densidade. Por isso este documento não tenta encaixar gases nos campos de peso já existentes: a reintrodução de \texttt{GASOSO} deve vir acompanhada desses campos e dessa fórmula alternativa, não apenas do valor no enum.

\subsection{Horários de Turma}
A entidade \textit{Horario\_Turma} (dia da semana, hora de início/fim e local por turma) foi removida da V1 por não possuir fluxo de uso definido nas telas nem nas regras de negócio. Em versões futuras, ela pode ser reintroduzida para permitir que o sistema exiba grades horárias, detecte conflitos de sala e envie lembretes automáticos antes do início das aulas.

\subsection{Lotação física de professores (V2)}
Q01: vinculação a prédio, gabinete e sala não integra a V1 por insuficiência de dados estruturados. Integração futura de dados institucionais de lotação e alocação de salas permitirá filtragem espacial; depende de disponibilização e validação dos dados pela UENF. Na V1, busca de professor limita-se a Nome, Matéria, Centro e Laboratório.

\subsection{Edição de Resumo e Especificação de Reagente (V2)}
Na V1, \textit{Resumo\_Reagente} e \textit{Especificacao\_Reagente} podem ser cadastrados, mas \textbf{não existe edição posterior} de Resumo ou Especificação pela interface, e não há endpoint genérico de edição desses objetos. Não pertencem à V1 \textit{Resumo.version}, \textit{EditSession} nem arquitetura equivalente de travamento otimista; a concorrência de edição desses objetos fica deliberadamente para a V2.

A V2 deverá estudar edição auditável de \textit{Resumo\_Reagente} e \textit{Especificacao\_Reagente}, incluindo, quando a implementação futura for projetada: autorização; concorrência entre gestores; versionamento; preservação de histórico; auditoria; e efeitos sobre projeções/cache. Nada disso deve ser implementado agora.

\subsection{Correção auditável da validade de Frasco\_Reagente (V2)}
Na V1 não existe operação de edição ou correção posterior da validade cadastrada de um frasco. Erro excepcional desse tipo não possui correção autônoma pela interface, e não há API genérica de correção. A V2 deverá estudar três modalidades futuras:

\begin{enumerate}
    \item \texttt{DESCONHECIDA\_PARA\_CONHECIDA} -- passar a conhecer uma validade antes desconhecida;
    \item \texttt{DATA\_INCORRETA\_PARA\_DATA\_CORRETA} -- corrigir uma data cadastrada incorretamente;
    \item \texttt{CONHECIDA\_PARA\_DESCONHECIDA} -- registrar que a validade antes considerada conhecida é, na verdade, desconhecida.
\end{enumerate}

A implementação futura deverá preservar o valor anterior no histórico; preservar autoria, instante, motivo e evidência; não oferecer edição genérica \texttt{campo/valor}; recalcular campos derivados no servidor; não confundir correção documental com ``revalidar validade''; e não reescrever fatos históricos de empréstimos anteriores. Essas três modalidades não são implementadas nesta rodada.

\subsection{Captura automática de pesagens e correção metrológica auditável (versões futuras)}
Na V1, toda pesagem é \textbf{entrada manual}; a interface usa barreiras preventivas de dupla digitação independente, unidade fixa em g, contexto visual, alertas de plausibilidade e revisão final antes do commit (Seção~8, UI-18). A V1 \textbf{não} possui correção administrativa metrológica após a persistência, nem editor genérico de peso, nem tipo específico \texttt{ERRO\_TRANSCRICAO\_PESO}: uma nova pesagem não reescreve automaticamente o fato histórico anterior. A decisão humana de M6 (HQ-M6-001) mantém o conjunto metrológico fechado em três rotas físicas tipadas.

\subsubsection{Estudo futuro A — integração direta com a balança}
Estudar a integração das balanças do LCQUI diretamente ao sistema, para reduzir ou eliminar a transcrição manual. A arquitetura futura deverá avaliar, sem implementar agora: interfaces realmente disponíveis nas balanças do LCQUI (USB, serial, rede ou outro protocolo do equipamento); captura automática da leitura; identificação do instrumento; \texttt{timestamp} confiável; unidade; resolução/casas decimais; estabilidade da leitura; associação inequívoca entre frasco, operação e leitura; trilha de auditoria; autenticação/autorização; falhas de comunicação; modo offline/fallback manual; validação da interface; preservação do dado bruto; e prevenção de duplicidade/replay. Não escolher tecnologia/protocolo sem inventário real das balanças.

\subsubsection{Estudo futuro B — correção metrológica auditável}
Estudar se versões futuras devem suportar correções metrológicas administrativas excepcionalíssimas quando existir evidência independente suficiente de erro de transcrição. Não assumir que serão implementadas. Se estudadas, deverão responder: quais tipos fechados de correção existem; qual evidência é obrigatória; quem pode realizar; necessidade de segunda aprovação; preservação integral do valor original; autoria; \texttt{timestamp}; motivo; documento/evidência vinculada; efeito em FLOW; efeito em saldos; impossibilidade de alterar fisicamente o passado; e proibição de editor genérico \texttt{campo/valor}. É um estudo de versão futura, não requisito obrigatório de V1/V2.
